\chapter[IMPLEMENTAÇÃO DO SISTEMA]{Implementação do Sistema}
\label{cap:implementacao}

Este capítulo descreve a implementação do sistema projetado no
Capítulo~\ref{cap:projeto}. São apresentados a organização do código, a
configuração, a coleta, o processamento, a persistência, a linha de comando, o
servidor e o cliente web, o módulo de avaliação, o teste de regressão e um
guia para conduzir uma nova análise. As
tecnologias citadas foram apresentadas na Seção~\ref{sec:tecnologias}.

% -----------------------------------------------------------------------------
\section{Organização do código}
\label{sec:organizacao}

O código foi escrito em Python e organizado em pacotes, cada um com uma
responsabilidade, conforme a Tabela~\ref{tab:pacotes}. As dependências seguem
um único sentido: o \textit{pipeline} usa os pacotes de configuração, domínio e
utilitários, o servidor web usa o \textit{pipeline} e o banco, e nenhum pacote
de nível inferior depende da interface.

\begin{table}[htbp]
\centering
\caption{Organização do código-fonte em pacotes}
\label{tab:pacotes}
\small
\begin{tabularx}{\textwidth}{>{\ttfamily}l X}
\toprule
\normalfont Pacote ou módulo & Responsabilidade \\
\midrule
config/                & Variáveis de ambiente e configuração das heurísticas. \\
domain/                & Entidade \texttt{Edital}, que também define a tabela
                         do banco. \\
database/, repositories/ & Conexão com o banco e operações de leitura e
                         gravação de editais. \\
pipeline/search/       & Geração, priorização e execução das consultas e
                         estatísticas de desempenho das consultas. \\
pipeline/enrichment/   & \textit{Download} das páginas, extração de texto e
                         do domínio. \\
pipeline/filtering/    & Filtros de listagem e de relevância e extração de
                         prazo. \\
pipeline/processing/   & Deduplicação por URL, assinatura de conteúdo e
                         similaridade de Jaccard. \\
pipeline/pipeline.py   & Orquestração da coleta e do processamento. \\
pipeline/snapshot.py   & Gravação e leitura dos \textit{snapshots}. \\
pipeline/evaluation/   & Métricas, avaliação, análise de sensibilidade e busca
                         em grade. \\
utils/                 & Normalização de URLs e de texto. \\
web/                   & Servidor Flask, \textit{template} HTML, estilos e
                         \textit{scripts} da interface. \\
main.py                & Ponto de entrada da linha de comando. \\
\bottomrule
\end{tabularx}
\fonte{Elaborado pelo autor.}
\end{table}

% -----------------------------------------------------------------------------
\section{Configuração}
\label{sec:impl_configuracao}

A configuração tem duas partes. A primeira reúne os dados do ambiente de
execução: a chave de acesso à API de busca, o endereço de conexão com o banco
e os parâmetros da atualização automática (se ela está ligada, o horário
diário da coleta e o fuso horário, por padrão meia-noite no horário de
Brasília). Esses valores são lidos de variáveis
de ambiente, carregadas de um arquivo local que não é versionado, o que mantém
as credenciais fora do código-fonte (RNF07).

A segunda parte é a configuração das heurísticas, implementada como um conjunto
de classes de dados imutáveis (\textit{frozen dataclasses}) do Python. Há uma
classe para cada grupo de parâmetros (listagem, relevância, prazo e
deduplicação), e uma classe principal, \texttt{HeuristicConfig}, que os
agrega. A instância \texttt{DEFAULT} contém a configuração calibrada
(Seção~\ref{sec:configuracao}) e é usada pelo \textit{pipeline} quando nenhuma
outra é informada. A função \texttt{with\_overrides} recebe uma configuração e
um dicionário de alterações, que associa o nome de cada parâmetro, como
\texttt{relevance.threshold}, ao novo valor, e devolve uma nova configuração com os valores trocados. É por essa função que a
análise de sensibilidade e a busca em grade geram as configurações avaliadas.
O método \texttt{to\_dict} converte a configuração em dicionário, para que seja
gravada no manifesto de cada coleta.

% -----------------------------------------------------------------------------
\section{Coleta}
\label{sec:impl_coleta}

\subsection{Geração e priorização das consultas}

A função \texttt{generate\_queries} gera as consultas em três grupos: consultas
por domínio, que combinam seis termos-base com quatro domínios institucionais
por meio do operador \texttt{site:} e acrescentam o ano de referência (o ano
corrente no horário de Brasília, obtido a cada geração de consultas);
consultas temáticas, formadas a partir de uma lista configurável de temas; e
duas consultas regionais. Após a remoção de repetições, as consultas com
\texttt{site:} são colocadas à frente e a lista é limitada às 20 primeiras.
Como há 24 consultas por domínio, as temáticas e as regionais ficam fora desse
limite na versão atual do código. Na coleta de calibração, feita com uma versão
anterior que tinha quatro termos-base, duas consultas temáticas chegaram a ser
executadas; na coleta de teste, todas as consultas executadas foram por
domínio.

As consultas restantes são ordenadas pela função \texttt{prioritize\_queries},
que atribui a cada uma a soma de uma pontuação fixa e de uma pontuação
histórica. A pontuação fixa vale 3 pontos para a presença do operador
\texttt{site:}, 2 para o ano de referência, 2 para o termo ``edital'' e 1 para
a expressão ``inscrições abertas''. A pontuação histórica é lida do arquivo de
estatísticas das consultas e é atualizada após cada busca pela função
\texttt{update\_query\_stats}: uma consulta sem resultados perde 2 pontos, e
cada resultado soma 2 pontos se o título contém ``edital'', 1 se o resumo
contém ``inscrição'' ou ``prazo'' e 1 se o resumo contém o ano. Desse modo, as
consultas que trouxeram resultados promissores no passado são executadas
primeiro nas coletas seguintes.

\subsection{Busca e \textit{download}}

A função \texttt{search\_api} envia cada consulta à API Tavily por uma
requisição HTTP POST, com profundidade de busca básica e no máximo 5
resultados por consulta, e devolve para cada resultado o título, a URL e o
resumo. A coleta é encerrada quando se acumulam 50 resultados brutos, o que,
na prática, corresponde a cerca de 10 consultas por coleta (RNF04). Erros de
conexão, de tempo de resposta ou de requisição são registrados e tratados como
ausência de resultados, de modo que uma falha da API em uma consulta não
interrompe a coleta (RNF06).

Os resultados são deduplicados pela forma normalizada da URL, e de cada URL é
gerado um identificador estável, formado pelos 16 primeiros caracteres
hexadecimais do \textit{hash} SHA-256 da URL normalizada. Esse identificador
liga o item do manifesto às linhas dos arquivos de predições e de rótulos.

A função \texttt{fetch\_raw} baixa cada página com tempo limite de 15 segundos
e guarda o conteúdo em \textit{bytes}, sem interpretá-lo, junto com o código de
\textit{status} HTTP, o tipo de conteúdo, a URL final após redirecionamentos e
a codificação de caracteres. Quando o \textit{download} falha, o erro é
registrado no item e a coleta prossegue. Nas coletas utilizadas neste
trabalho, o \textit{download} falhou em 7 dos 47 itens da coleta de calibração
e em 13 dos 45 itens da coleta de teste, por erros de conexão e por respostas
HTTP 429 (excesso de requisições) e 404 (página inexistente). Nesses casos, o
processamento usa apenas o resumo devolvido pela API.

\subsection{Gravação do \textit{snapshot}}

A função \texttt{save\_snapshot} grava o conteúdo de cada página no
subdiretório \texttt{raw/}, com nome igual ao \textit{hash} SHA-256 do conteúdo
e extensão inferida do tipo (\texttt{.html}, \texttt{.pdf} ou \texttt{.bin}).
Páginas com conteúdo idêntico compartilham o mesmo arquivo. Em seguida, grava
o manifesto em JSON, com o instante da coleta, o \textit{commit} corrente,
obtido do Git, as versões das bibliotecas de extração, as consultas
executadas, a configuração e os metadados dos itens. A coleta devolve o
\textit{snapshot} lido do disco, de modo que o processamento de uma coleta
nova e o reprocessamento de uma coleta antiga seguem exatamente o mesmo
caminho no código.

% -----------------------------------------------------------------------------
\section{Processamento}
\label{sec:impl_processamento}

O processamento é implementado pelo método \texttt{process} da classe
\texttt{EditalPipeline}, que percorre os itens do \textit{snapshot} e aplica a
cadeia de etapas da Seção~\ref{sec:arquitetura}. Para cada item, o texto
extraído da página é comparado com o resumo da API, e o mais longo é usado nas
heurísticas. A extração trata separadamente os documentos PDF, reconhecidos
pela assinatura \texttt{\%PDF} no início do conteúdo ou pelo tipo declarado, e
os documentos HTML, dos quais são removidos \textit{scripts} e estilos antes da
extração do texto.

Quando a persistência está ativa, o item cuja URL já existe no banco é
descartado antes dos filtros, com o motivo \texttt{url\_duplicada}. Essa
verificação evita reprocessar editais já gravados por coletas anteriores e
não ocorre no modo de reprocessamento.

Os filtros são implementados em módulos próprios. O filtro de listagem é uma
função que calcula a pontuação $L$ (Equação~\ref{eq:listagem}) e a compara ao
limiar. O filtro de relevância é uma classe que compila, na inicialização,
uma expressão regular para cada termo do vocabulário, de modo que só
ocorrências de palavras inteiras sejam contadas. A extração de prazo é uma
classe que recebe na inicialização a data de referência, lida do manifesto, e
executa os passos da Seção~\ref{subsec:prazo} com as opções ligadas na
configuração calibrada. Quando nenhum prazo é encontrado no texto completo,
a extração é repetida sobre os primeiros 2.000 caracteres.

A deduplicação por conteúdo mantém, durante o processamento de uma coleta, o
conjunto das assinaturas MD5 já vistas e a lista dos textos já aceitos. Um item
é descartado se sua assinatura já consta do conjunto ou se a similaridade de
Jaccard com algum texto aceito supera o limiar. Os itens que passam por todas
as etapas são convertidos em objetos \texttt{Edital}, com o domínio extraído
da URL, e gravados no banco.

Ao final, o método grava o arquivo de predições, com uma linha por item. Cada
linha contém o identificador, a URL e o título do item, a origem do texto
(página ou resumo), as decisões e pontuações de cada etapa, a frase de onde o
prazo foi extraído e o motivo da decisão final, que é um dos seguintes:
\texttt{url\_duplicada}, \texttt{listing\_page}, \texttt{not\_relevant},
\texttt{no\_deadline}, \texttt{expired\_deadline},
\texttt{duplicate\_signature}, \texttt{duplicate\_similarity} ou
\texttt{accepted}.

% -----------------------------------------------------------------------------
\section{Persistência}
\label{sec:impl_persistencia}

A entidade \texttt{Edital} é declarada como uma classe do SQLModel que define,
ao mesmo tempo, o objeto manipulado pelo \textit{pipeline} e a tabela do
banco. A tabela é criada, se ainda não existir, pela função \texttt{create\_db},
e pode ser esvaziada pela função \texttt{clear\_database}. O acesso aos dados é
concentrado na classe \texttt{EditalRepository}, que oferece as operações de
gravar um edital, verificar se uma URL já existe e listar todos os editais.
Cada operação abre e fecha a própria sessão com o banco.

A criação automática da tabela não altera tabelas já existentes. Em um banco
criado com uma versão anterior do modelo, a coluna de prazo foi criada com o
tipo data e hora, em vez de data. O servidor web trata esse caso ao converter
o valor lido para data antes de enviá-lo à interface.

% -----------------------------------------------------------------------------
\section{Linha de comando}
\label{sec:impl_cli}

O \textit{pipeline} pode ser executado diretamente pelo módulo
\texttt{main.py}. Sem argumentos, executa uma coleta completa, com gravação no
banco. A opção \texttt{-{}-reset} esvazia o banco antes da coleta, e a opção
\texttt{-{}-replay}, seguida do diretório de uma coleta, reprocessa o
\textit{snapshot} correspondente sem acessar a Web nem o banco, gravando
apenas o arquivo de predições. Ao final, os editais aceitos são listados no
terminal, com título, URL, domínio, prazo, frase de origem e pontuação.

% -----------------------------------------------------------------------------
\section{Servidor web}
\label{sec:impl_servidor}

O servidor é uma aplicação Flask com duas rotas. A rota \texttt{/} devolve a
página da interface, gerada a partir de um \textit{template} Jinja2. A rota
\texttt{/api/editais} consulta o banco e devolve, em JSON, os editais com
prazo igual ou posterior à data corrente, ordenados do prazo mais próximo ao
mais distante e, em caso de empate, da maior para a menor pontuação de
relevância. Antes do envio, são aplicados os dois filtros de exibição da
Seção~\ref{sec:decisoes}. Para cada edital, a resposta contém o
identificador, o título, a URL, o domínio e o prazo no formato AAAA-MM-DD, e,
para a lista, o total de editais, o instante da última coleta, o instante
previsto para a próxima e a indicação de que uma coleta está em andamento.
Se o banco estiver indisponível, a rota responde com o código HTTP 503 e uma
mensagem de erro (RNF06).

A atualização automática é implementada pela classe \texttt{Updater}, cuja
rotina principal é executada em uma \textit{thread} de segundo plano, iniciada
junto com o servidor. O instante da última coleta é lido do manifesto mais
recente do diretório de coletas, e não da memória do processo. Assim,
reiniciar o servidor não antecipa nem adia a coleta. A próxima coleta é
agendada para a primeira meia-noite, no fuso \texttt{America/Sao\_Paulo},
posterior à última coleta; se o servidor estava desligado nesse horário, a
coleta perdida é executada assim que ele voltar. A
rotina calcula o tempo até a próxima coleta e aguarda no máximo uma hora antes
de reavaliar o horário. Uma coleta que falha é registrada e repetida após uma
hora. Cada coleta automática grava um \textit{snapshot} completo, como as
coletas feitas pela linha de comando, o que forma um acervo de coletas
disponível para rotulação e recalibração (Seção~\ref{sec:impl_recalibracao}).
Um bloqueio (\textit{lock}) impede que duas coletas sejam executadas ao
mesmo tempo, e, pelo mesmo motivo, o recarregador automático do Flask, que
iniciaria um segundo processo, foi desativado.

Por fim, o servidor registra explicitamente os tipos de conteúdo dos arquivos
JavaScript e CSS. Em alguns ambientes Windows, o registro do sistema associa a
extensão \texttt{.js} a texto simples, e o navegador recusa carregar módulos
JavaScript servidos com esse tipo.

% -----------------------------------------------------------------------------
\section{Cliente web}
\label{sec:impl_cliente}

O cliente separa estrutura, apresentação e comportamento em arquivos
distintos. O \textit{template} HTML contém apenas a estrutura da página, os
estilos estão em três arquivos CSS (base, \textit{layout} e componentes) e o
comportamento está em seis módulos JavaScript, descritos na
Tabela~\ref{tab:modulos_js}. Somente o módulo principal é carregado pela
página, e ele importa os demais.

\begin{table}[htbp]
\centering
\caption{Módulos JavaScript da interface}
\label{tab:modulos_js}
\small
\begin{tabularx}{\textwidth}{>{\ttfamily}l X}
\toprule
\normalfont Módulo & Responsabilidade \\
\midrule
main.js    & Ponto de entrada: exibe a lista guardada e liga o botão
             ``Atualizar''. \\
api.js     & Requisição à rota \texttt{/api/editais}. \\
storage.js & Leitura e gravação da lista no \textit{localStorage}. \\
render.js  & Montagem dos cards, da contagem e da linha de \textit{status}. \\
toast.js   & Mensagens temporárias de sucesso e de erro. \\
utils.js   & Proteção do HTML, validação de \textit{links} e formatação de
             datas. \\
\bottomrule
\end{tabularx}
\fonte{Elaborado pelo autor.}
\end{table}

Ao carregar, a página lê a lista guardada no navegador e a exibe, sem
requisição ao servidor. Ao pressionar ``Atualizar'', o botão é desabilitado, a
lista é solicitada ao servidor e, em caso de sucesso, é guardada junto com o
instante do carregamento e redesenhada. Uma mensagem temporária informa o
número de editais recebidos ou o erro ocorrido. Se o servidor estiver
executando uma coleta naquele momento, a mensagem sugere uma nova tentativa
em alguns minutos. As falhas de acesso ao armazenamento local, que ocorrem,
por exemplo, em navegação privada, são ignoradas, e a página continua
funcionando sem o \textit{cache}.

Como títulos e URLs vêm de páginas da Web, todo texto é convertido em entidades
HTML antes de ser inserido na página, e só são aceitos \textit{links} com os
protocolos HTTP e HTTPS. Isso impede que um título ou uma URL maliciosos
executem código no navegador do usuário (RNF07). Os \textit{links} abrem em
nova aba, sem dar à página de destino acesso à página de origem.

A grade de cards usa o sistema de colunas do Bootstrap, com os pontos de quebra
descritos na Seção~\ref{sec:proj_interface}. Os estilos foram escritos
partindo da menor tela para a maior: as regras gerais valem para celulares, e
as regras para telas maiores são acrescentadas por consultas de mídia. Na
largura de celular, o botão ``Atualizar'' ocupa a largura do quadro, e todos os
elementos de toque têm ao menos 44 \textit{pixels} de altura. Títulos longos
são limitados a três linhas, e domínios longos são abreviados com reticências,
para que os cards de uma mesma linha mantenham a mesma estrutura.

% -----------------------------------------------------------------------------
\section{Módulo de avaliação}
\label{sec:impl_avaliacao}

O módulo de avaliação implementa o procedimento da
Seção~\ref{sec:calibracao} em cinco componentes, todos executados pela linha de
comando sobre o diretório de uma coleta:

\begin{itemize}
  \item \textbf{Métricas}: funções que calculam a matriz de confusão, a
    precisão, a revocação e as medidas $F_{0{,}5}$, $F_1$ e $F_2$
    (Equação~\ref{eq:fbeta}).
  \item \textbf{Avaliador}: compara o arquivo de predições com o arquivo de
    rótulos e grava as métricas de cada etapa. Na primeira execução sobre uma
    coleta, cria o arquivo de rótulos em branco, com uma linha por item, para
    a rotulação manual, podendo omitir os itens já rotulados em outras
    coletas.
  \item \textbf{Experimentos}: calcula as pontuações das três heurísticas para
    todos os itens, independentemente da cadeia, e reproduz a cadeia a partir
    dessas pontuações. Implementa as duas formas de avaliação da
    Seção~\ref{subsec:modos}. Como o texto de cada item é extraído uma única
    vez e reutilizado, cada nova configuração é avaliada sem novo acesso aos
    arquivos da coleta.
  \item \textbf{Análise de sensibilidade}: varia cada parâmetro ao longo de
    sua grade de valores, grava uma tabela e um gráfico por parâmetro, gerado
    com a biblioteca Matplotlib \cite{hunter2007}, e uma tabela-resumo com a
    amplitude e o valor sugerido de cada parâmetro. Com uma opção adicional,
    compara as variantes estruturais com a configuração de base.
  \item \textbf{Busca em grade}: avalia as combinações da
    Tabela~\ref{tab:grade}, ordena-as pelo critério lexicográfico e executa a
    validação cruzada e o \textit{bootstrap} descritos na
    Seção~\ref{subsec:etapaD}, com sementes fixas para que os resultados sejam
    reproduzíveis.
\end{itemize}

% -----------------------------------------------------------------------------
\section{Teste de regressão}
\label{sec:impl_regressao}

O teste de regressão descrito na Seção~\ref{sec:coleta} é executado em dois
comandos. O primeiro reprocessa o \textit{snapshot} de calibração no modo sem
persistência e grava as predições em um arquivo temporário, para não
sobrescrever o arquivo de referência. O segundo compara os dois arquivos
byte a byte e informa se são idênticos. A comparação é feita sobre o
conteúdo, e não apenas sobre o tamanho e a data dos arquivos.

Como o arquivo de predições registra todos os itens, com as decisões e
pontuações de cada etapa e o motivo da decisão final, qualquer alteração de
comportamento, mesmo em um item descartado, produz uma diferença. O teste não
verifica se as decisões estão corretas, papel do módulo de avaliação, mas
apenas se mudaram. Ele foi executado após cada refatoração, e o arquivo de
referência só foi substituído no congelamento da configuração
(Seção~\ref{subsec:etapaE}).

Como o arquivo de referência é versionado, o teste também pode ser feito
reprocessando o \textit{snapshot} no próprio diretório e pedindo ao Git as
diferenças do arquivo de predições. Uma saída vazia indica identidade:

\begin{footnotesize}
\begin{verbatim}
> python main.py --replay datasets/run_2026-10-27T17-50-29
> git diff --stat -- datasets/run_2026-10-27T17-50-29/predictions.csv
>
\end{verbatim}
\end{footnotesize}

% -----------------------------------------------------------------------------
\section{Guia para uma nova análise}
\label{sec:impl_recalibracao}

As páginas das instituições de fomento mudam de estrutura, o vocabulário dos
editais varia e novas fontes podem ser incluídas nas consultas. Por isso, a
configuração calibrada neste trabalho não deve ser considerada definitiva.
Esta seção é um guia para conduzir uma nova análise do zero, com coletas e
rótulos novos, seguindo as Etapas A a E da Seção~\ref{sec:calibracao}. Para
cada passo são dados o objetivo, os comandos e o retorno esperado.

Todos os comandos são executados a partir da raiz do projeto, em um terminal
PowerShell; nos comandos longos, o acento grave (\texttt{`}) ao final da linha
indica continuação. \texttt{run\_CAL} e \texttt{run\_TESTE} representam os
diretórios das coletas de calibração e de teste da nova análise, nomeados pelo
instante da coleta (por exemplo, \texttt{run\_2027-03-02T00-00-05}). Os
retornos mostrados são exemplos, obtidos com as coletas deste trabalho; numa
nova análise, os números serão outros, mas o formato é o mesmo. Trechos
repetitivos foram omitidos (\texttt{[...]}).

O guia não exige alteração do código para ser seguido em outro ano: o ano das
consultas é o ano corrente (Seção~\ref{sec:impl_coleta}), e a data de
referência de cada coleta é lida do seu manifesto. As alterações de código
previstas são apenas as que a própria análise decidir fazer, nos Passos 6 e 8.

\begin{table}[htbp]
\centering
\caption{Passos de uma nova análise}
\label{tab:guia}
\small
\begin{tabularx}{\textwidth}{c l X}
\toprule
Passo & Etapa & Ação \\
\midrule
1 & --- & Preparar o ambiente e confirmar o estado do código. \\
2 & --- & Fazer as coletas de calibração e de teste. \\
3 & B   & Rotular as duas coletas; reservar a de teste. \\
4 & A   & Avaliar a configuração vigente na calibração. \\
5 & C   & Executar a análise de sensibilidade. \\
6 & D   & Implementar e comparar variantes estruturais. \\
7 & D   & Executar a busca em grade e a validação. \\
8 & E   & Congelar a nova configuração. \\
9 & E   & Avaliar a coleta de teste, uma única vez. \\
10 & --- & Remover mecanismos sem efeito e registrar a análise. \\
\bottomrule
\end{tabularx}
\fonte{Elaborado pelo autor.}
\end{table}

\subsection{Passo 1: preparação do ambiente}

Instalam-se as dependências nas versões fixadas e inicia-se o contêiner do
banco. O arquivo \texttt{.env}, na raiz do projeto, deve conter a chave da API
de busca (\texttt{API\_KEY}) e o endereço do banco (\texttt{DATABASE\_URL}).

\begin{footnotesize}
\begin{verbatim}
> pip install -r requirements.txt
> docker start postgres-editais
\end{verbatim}
\end{footnotesize}

Em seguida, executa-se o teste de regressão (Seção~\ref{sec:impl_regressao})
sobre a coleta de calibração da análise anterior, que está versionada com o seu
arquivo de predições de referência. Uma saída vazia confirma que o código se
comporta exatamente como no último congelamento; uma diferença indica uma
alteração não registrada, que deve ser esclarecida antes de prosseguir, pois a
nova análise partiria de uma configuração desconhecida.

\begin{footnotesize}
\begin{verbatim}
> python main.py --replay datasets/run_2026-10-27T17-50-29
> git diff --stat -- datasets/run_2026-10-27T17-50-29/predictions.csv
>
\end{verbatim}
\end{footnotesize}

\subsection{Passo 2: coletas}

São necessárias duas coletas, uma para a calibração e outra para o teste,
feitas em datas distintas para que o teste represente dados não vistos. Há
duas formas de obtê-las:

\begin{itemize}
  \item \textbf{Coletas do servidor}. Com o servidor web em execução, uma
    coleta é feita todos os dias à meia-noite (Seção~\ref{sec:impl_servidor}),
    e cada uma grava um \textit{snapshot} em \texttt{datasets/}. Basta escolher
    dois desses diretórios, de datas afastadas.
  \item \textbf{Linha de comando}. Cada execução faz uma coleta completa:
\end{itemize}

\begin{footnotesize}
\begin{verbatim}
> python main.py
Mantendo dados existentes...
Resultados brutos: <n>
Após deduplicação: <n>
Snapshot salvo em: datasets\run_AAAA-MM-DDTHH-MM-SS
[...]
Avaliação salva em: datasets\run_AAAA-MM-DDTHH-MM-SS\predictions.csv
Editais válidos: <aceitos>/<itens>
\end{verbatim}
\end{footnotesize}

Nas duas formas, a coleta é feita com persistência, e os itens cujas URLs já
estão no banco são descartados com o motivo \texttt{url\_duplicada}
(Seção~\ref{sec:impl_processamento}). As predições gravadas na coleta,
portanto, não servem para a avaliação. Elas são preservadas com outro nome, e
o \textit{snapshot} é reprocessado sem banco, o que gera predições para todos
os itens. O procedimento é repetido para a coleta de teste:

\begin{footnotesize}
\begin{verbatim}
> Rename-Item datasets/run_CAL/predictions.csv predictions_coleta.csv
> python main.py --replay datasets/run_CAL
[PROCESSANDO]
[URL]: http://resultado.cnpq.br/5170965351053941
[SCORE]: 91
[EXPIRED DEADLINE]
Avaliação salva em: datasets\run_CAL\predictions.csv
Editais válidos: 2/47
Editais encontrados: 2
\end{verbatim}
\end{footnotesize}

\subsection{Passo 3: rotulação e reserva do teste}

Na primeira execução do avaliador sobre uma coleta, é criado o arquivo de
rótulos em branco, com uma linha por item:

\begin{footnotesize}
\begin{verbatim}
> python -m pipeline.evaluation.evaluator datasets/run_CAL
labels.csv criado em datasets\run_CAL\labels.csv. Preencha os rótulos e rode
novamente.
\end{verbatim}
\end{footnotesize}

Na coleta de teste, nenhum item usado na calibração pode ser usado na
avaliação. A opção \texttt{-{}-excluir}, seguida dos diretórios da coleta de
calibração e de qualquer outro conjunto rotulado anterior, cria o arquivo de
rótulos sem os itens que já constam dos rótulos dessas coletas. Como o
identificador de cada item é derivado da URL normalizada
(Seção~\ref{sec:impl_coleta}), um mesmo endereço recebe o mesmo identificador
em coletas diferentes. A opção só atua na criação do arquivo e nunca altera
rótulos já preenchidos:

\begin{footnotesize}
\begin{verbatim}
> python -m pipeline.evaluation.evaluator datasets/run_TESTE `
    --excluir datasets/run_CAL
19 de 45 itens já rotulados em datasets/run_CAL foram omitidos.
labels.csv criado em datasets\run_TESTE\labels.csv. Preencha os rótulos e rode
novamente.
\end{verbatim}
\end{footnotesize}

Os rótulos da Tabela~\ref{tab:rotulos} são preenchidos com \texttt{yes} ou
\texttt{no}, e a data máxima de inscrição no formato AAAA-MM-DD, segundo o guia
de rotulação (\texttt{datasets/GUIA\_ROTULACAO.md}) e suas duas regras de
consistência (Seção~\ref{sec:dados}). As duas coletas devem ser rotuladas com
os mesmos critérios. Se a planilha for editada no Excel, deve ser salva no
formato CSV UTF-8 e fechada antes dos comandos seguintes, pois o Excel bloqueia
a gravação dos arquivos abertos.

A partir deste passo, a coleta de teste fica reservada: nenhum comando do
módulo de avaliação é executado sobre ela até o Passo 9. Admite-se apenas a
conferência manual do formato e da consistência dos rótulos.

\subsection{Passo 4: avaliação da configuração vigente}

Antes de qualquer ajuste, mede-se o desempenho da configuração vigente na
coleta de calibração. O avaliador grava as métricas da execução em cadeia em
\texttt{metrics.csv}, e o módulo de experimentos grava as pontuações das
heurísticas em \texttt{features.csv} e imprime as métricas por etapa e
operacionais (Seções~\ref{subsec:operacionais} e~\ref{subsec:modos}). Esses
números são a referência para as decisões dos passos seguintes, e o
\texttt{metrics.csv} deve ser preservado com outro nome antes do Passo 8, pois
será regravado com a nova configuração.

\begin{footnotesize}
\begin{verbatim}
> python -m pipeline.evaluation.evaluator datasets/run_CAL
listing      n=47  P=0.750 R=0.333 F1=0.462 F2=0.375
relevance    n=47  P=0.667 R=0.889 F1=0.762 F2=0.833
is_edital    n=47  P=0.000 R=0.000 F1=0.000 F2=0.000
end_to_end   n=47  P=0.000 R=0.000 F1=0.000 F2=0.000
deadline     n=47  P=0.000 R=0.000 F1=0.000 F2=0.000

> python -m pipeline.evaluation.experiments datasets/run_CAL
listing            n=47  P=0.750 R=0.333 F0.5=0.600 F1=0.462 F2=0.375
relevance          n=47  P=0.643 R=1.000 F0.5=0.692 F1=0.783 F2=0.900
deadline_exists    n=20  P=0.947 R=0.947 F0.5=0.947 F1=0.947 F2=0.947
deadline_open      n=19  P=0.000 R=0.000 F0.5=0.000 F1=0.000 F2=0.000
deadline_date      n=19  acertos=10 acc=0.526
editais_perdidos   n=20  perdidos=1
end_to_end         n=47  P=0.000 R=0.000 F0.5=0.000 F1=0.000 F2=0.000
aceitos_data_certa n=2   acertos=0 acc=0.000
\end{verbatim}
\end{footnotesize}

As métricas de relevância diferem entre os dois comandos porque o avaliador lê
as decisões da execução em cadeia, na qual um item descartado como listagem não
chega a ter a relevância decidida e é contado como não relevante, enquanto o
módulo de experimentos avalia cada heurística sobre todos os itens. As
decisões da análise apoiam-se nas métricas do módulo de experimentos,
sobretudo em \texttt{editais\_perdidos} e nos itens errados aceitos (os falsos
positivos de \texttt{end\_to\_end}).

\subsection{Passo 5: análise de sensibilidade}

A análise OAT (Etapa C) varia cada parâmetro numérico ao longo de uma grade de
valores, mantidos os demais na configuração de base, e leva cerca de dois
minutos sobre uma coleta de 47 itens. Para cada parâmetro, imprime a amplitude
da métrica-alvo e do $F_1$ do resultado final, o valor atual, o valor sugerido e
os editais perdidos e itens errados aceitos nos dois valores:

\begin{footnotesize}
\begin{verbatim}
> python -m pipeline.evaluation.sensitivity_analysis datasets/run_CAL
listing.threshold          amp=0.682 amp_e2e=0.000 atual=6 sugerido=6
                           perdidos 1->1 errados 2->2
listing.date_high          amp=0.099 amp_e2e=0.000 atual=20 sugerido=30
                           perdidos 1->1 errados 2->2
[...]
deadline.year_window       amp=0.000 amp_e2e=0.000 atual=10 sugerido=10
                           perdidos 1->1 errados 2->2
[...]
\end{verbatim}
\end{footnotesize}

Cada linha foi quebrada aqui para caber na página; no terminal, ocupa uma só.
Os resultados são gravados no subdiretório \texttt{sensitivity\_base/} da
coleta: uma tabela e um gráfico por parâmetro e a tabela-resumo
\texttt{impact.csv}, ordenada pela amplitude. Para analisar apenas alguns
parâmetros, seus nomes são passados após o diretório, e a tabela-resumo passa a
se chamar \texttt{impact\_<grupos>.csv}:

\begin{footnotesize}
\begin{verbatim}
> python -m pipeline.evaluation.sensitivity_analysis datasets/run_CAL `
    listing.threshold deadline.year_window
\end{verbatim}
\end{footnotesize}

A leitura segue três regras:

\begin{itemize}
  \item um valor sugerido diferente do atual só é candidato a alteração se
    melhora o critério lexicográfico (Seção~\ref{subsec:etapaC}): não aumenta
    os itens errados, perde menos editais ou, empatados esses, melhora a
    métrica-alvo;
  \item quando a amplitude é nula, todos os valores empatam e o valor sugerido
    é apenas o centro da grade, sem significado. O parâmetro deve ser
    investigado no código, pois pode haver um defeito que o impede de atuar,
    como os encontrados neste trabalho (Seção~\ref{sec:res_sensibilidade});
  \item as grades de valores ficam no dicionário \texttt{GRID} do módulo e
    podem ser ampliadas quando o melhor valor estiver na borda da grade.
\end{itemize}

A configuração de base é a constante \texttt{BASE} do módulo, que deve ser
igual à configuração padrão vigente. Ela coincide hoje com a configuração
padrão; se uma análise anterior a tiver alterado, deve ser corrigida antes deste
passo.

\subsection{Passo 6: variantes estruturais}

As hipóteses de mudança de regra levantadas no passo anterior, como uma nova
forma de extrair a data, são implementadas como opções da configuração,
desligadas por padrão, no arquivo \texttt{config/heuristics.py} e no módulo do
filtro correspondente. Depois de cada implementação, o teste de regressão do
Passo 1 deve resultar em identidade, o que confirma que a opção desligada não
altera o comportamento. A variante é então cadastrada no dicionário
\texttt{VARIANTS} do módulo de análise de sensibilidade, com um nome e as
alterações de configuração que a ligam, e comparada à base:

\begin{footnotesize}
\begin{verbatim}
> python -m pipeline.evaluation.sensitivity_analysis datasets/run_CAL --variants
base            perdidos=1 errados=2 listing.f05=0.600 relevance.f2=0.900
                data=0.526
relevance_tema  perdidos=1 errados=2 listing.f05=0.600 relevance.f2=0.894
                data=0.526
\end{verbatim}
\end{footnotesize}

A tabela correspondente é gravada em \texttt{variants/variants.csv}. Uma
variante é aceita somente se atende ao critério lexicográfico; no exemplo, a
variante de vocabulário temático não muda editais perdidos nem itens errados e
reduz a métrica-alvo, e seria rejeitada. As variantes aceitas são ligadas na
constante \texttt{BASE}, e a análise do Passo 5 é repetida sobre a nova base.

Ao final deste passo, o código das variantes é registrado em um
\textit{commit}, com todas as opções ainda desligadas na configuração padrão.
O arquivo \texttt{config/heuristics.py} desse \textit{commit} contém, portanto,
a configuração anterior completa, inclusive as opções novas, e será usado no
Passo 9. Se nenhuma variante for implementada, anota-se o \textit{commit}
vigente:

\begin{footnotesize}
\begin{verbatim}
> git add config pipeline
> git commit -m "Variantes da nova análise, desligadas por padrão"
> git log --oneline -1
<commit_anterior> Variantes da nova análise, desligadas por padrão
\end{verbatim}
\end{footnotesize}

\subsection{Passo 7: busca em grade e validação}

Como a análise OAT não captura interações entre parâmetros, a busca em grade
(Etapa D) avalia todas as combinações dos parâmetros do filtro de listagem
definidas no dicionário \texttt{GRID} do módulo (Tabela~\ref{tab:grade}), a
partir da configuração \texttt{BASE}, e em seguida executa a validação cruzada
e o \textit{bootstrap}. Leva cerca de 15 segundos:

\begin{footnotesize}
\begin{verbatim}
> python -m pipeline.evaluation.grid_search datasets/run_CAL
melhor da grade: {'listing.url_pattern_weight': 0, 'listing.url_keyword_weight': 2,
  'listing.date_weight': 2, 'listing.date_high': 20,
  'listing.plural_editais_weight': 2, 'listing.threshold': 6, 'changed': 0,
  'lost': 1, 'wrong': 2, 'f05': 0.6000000000000001}
empatados no topo: 4
CV base       perdidos=10 errados=20
CV escolhida  perdidos=22 errados=20
estabilidade: [('0|2|2|20|2|6', 40), ('0|2|2|20|2|5', 8), ('0|2|3|20|2|5', 1)]
  de 50 dobras
bootstrap IC95% (escolhida - base): {'lost': (0, 0), 'wrong': (0, 0)}
\end{verbatim}
\end{footnotesize}

As linhas do retorno significam:

\begin{itemize}
  \item \texttt{melhor da grade}: a primeira combinação pelo critério
    lexicográfico, com o número de parâmetros que ela altera em relação à base
    (\texttt{changed}), os editais perdidos (\texttt{lost}), os itens errados
    aceitos (\texttt{wrong}) e a $F_{0{,}5}$ da listagem;
  \item \texttt{empatados no topo}: quantas combinações empatam com a
    primeira; entre elas, a ordenação já privilegia a que altera menos
    parâmetros;
  \item \texttt{CV}: soma, sobre as 50 partições da validação cruzada, dos
    editais perdidos e itens errados nos itens retidos, com a base e com a
    combinação escolhida usando apenas os itens de treino de cada partição;
  \item \texttt{estabilidade}: as combinações mais escolhidas nas partições;
  \item \texttt{bootstrap}: intervalo de confiança de 95\% da diferença entre
    a melhor combinação e a base, em editais perdidos e itens errados.
\end{itemize}

As tabelas completas são gravadas em \texttt{grid/grid\_listing.csv} e
\texttt{grid/cv\_listing.csv}. A melhor combinação só deve ser adotada se
\texttt{changed} for maior que zero, se a linha \texttt{CV} da combinação
escolhida não for pior que a da base e se o intervalo de confiança não indicar
piora. No exemplo, executado sobre a configuração já calibrada, a melhor
combinação é a própria base (\texttt{changed}~=~0), e não há o que alterar.

\subsection{Passo 8: congelamento}

Os valores e as opções escolhidos nos Passos 5 a 7 tornam-se o padrão das
classes de configuração em \texttt{config/heuristics.py}, e a constante
\texttt{BASE} é ajustada para coincidir com a nova configuração padrão. Neste
passo, apenas valores e opções mudam; nenhum mecanismo é removido, para que o
arquivo de configuração do \textit{commit} anotado no Passo 6 continue
executável pelo código atual no Passo 9.

Em seguida, o \textit{snapshot} de calibração é reprocessado, o que gera o novo
arquivo de referência do teste de regressão, as métricas da calibração são
regravadas, e tudo é registrado em um \textit{commit}:

\begin{footnotesize}
\begin{verbatim}
> Rename-Item datasets/run_CAL/metrics.csv metrics_pre_calibracao.csv
> python main.py --replay datasets/run_CAL
> python -m pipeline.evaluation.evaluator datasets/run_CAL
> python -m pipeline.evaluation.experiments datasets/run_CAL
> git add config/heuristics.py pipeline/evaluation datasets/run_CAL
> git commit -m "Congela configuração da nova análise"
\end{verbatim}
\end{footnotesize}

\subsection{Passo 9: avaliação no conjunto de teste}

Só depois do congelamento a coleta de teste é avaliada, uma única vez. Primeiro
com a configuração anterior, recuperando temporariamente o arquivo de
configuração do \textit{commit} anotado no Passo 6, e depois com a nova. Como ambas foram fixadas antes de
qualquer consulta ao teste, a comparação entre elas é legítima. O aviso do
avaliador corresponde aos itens omitidos no Passo 3.

\begin{footnotesize}
\begin{verbatim}
> python main.py --replay datasets/run_TESTE

> git checkout <commit_anterior> -- config/heuristics.py
> python -m pipeline.evaluation.experiments datasets/run_TESTE
[métricas da configuração anterior, no formato do Passo 4]
> git checkout HEAD -- config/heuristics.py

> python -m pipeline.evaluation.evaluator datasets/run_TESTE
[AVISO] 19 itens sem linha em labels.csv
listing      n=26  P=0.500 R=0.200 F1=0.286 F2=0.227
[...]
> python -m pipeline.evaluation.experiments datasets/run_TESTE
listing            n=26  P=0.500 R=0.200 F0.5=0.385 F1=0.286 F2=0.227
relevance          n=26  P=0.400 R=0.800 F0.5=0.444 F1=0.533 F2=0.667
deadline_exists    n=14  P=0.857 R=1.000 F0.5=0.882 F1=0.923 F2=0.968
deadline_open      n=14  P=0.667 R=0.667 F0.5=0.667 F1=0.667 F2=0.667
deadline_date      n=12  acertos=4 acc=0.333
editais_perdidos   n=14  perdidos=1
end_to_end         n=26  P=0.333 R=0.500 F0.5=0.357 F1=0.400 F2=0.455
aceitos_data_certa n=3   acertos=1 acc=0.333
\end{verbatim}
\end{footnotesize}

A segunda execução do módulo de experimentos deve ser a última, pois cada
execução regrava o \texttt{features.csv} com a configuração em uso, e o
arquivo que fica na coleta deve ser o da configuração nova. As duas saídas do
módulo de experimentos, anterior e nova, são registradas como resultado da
análise, nos moldes da Tabela~\ref{tab:resultados}.

Nenhum ajuste é feito depois desta avaliação. Uma nova alteração exige uma
nova análise, com nova coleta de teste, pois a atual deixa de ser independente
assim que suas métricas são conhecidas.

\subsection{Passo 10: limpeza e registro}

Por fim, os mecanismos que a análise de sensibilidade mostrou sem efeito podem
ser removidos do código, desde que o teste de regressão sobre a coleta de
calibração continue resultando em identidade. Os resultados do teste são
registrados em um \textit{commit} próprio:

\begin{footnotesize}
\begin{verbatim}
> python main.py --replay datasets/run_CAL
> git diff --stat -- datasets/run_CAL/predictions.csv
>
> git add datasets/run_TESTE <arquivos da limpeza>
> git commit -m "Avalia a nova configuração no conjunto de teste"
\end{verbatim}
\end{footnotesize}

\subsection{Cuidados}

\begin{itemize}
  \item \textbf{Saídas sobrescritas}. Os módulos gravam seus resultados em
    nomes fixos dentro do diretório da coleta. Executá-los sobre coletas de
    análises anteriores substitui os arquivos que fundamentam aqueles
    resultados; para consultá-las, as execuções devem ser feitas sobre uma
    cópia do diretório.
  \item \textbf{Histórico das consultas}. A pontuação histórica das consultas é
    guardada pelo texto exato de cada consulta. Como o texto inclui o ano, as
    consultas de um novo ano começam sem histórico.
  \item \textbf{Padrão de URL com ano}. A lista de padrões de URL do filtro de
    listagem contém \texttt{/editais-2026}. Com o peso $w_p$ igual a 0, os
    padrões não influenciam a pontuação; se a nova análise voltar a dar peso a
    eles, esse padrão deve ser atualizado.
  \item \textbf{Tamanho da amostra}. Com poucos editais abertos, uma diferença
    de um item altera muito as métricas do resultado final
    (Seção~\ref{sec:ameacas}). Entre as coletas disponíveis, convém escolher
    para a calibração e para o teste as que tenham mais editais abertos.
\end{itemize}
